\begin{lemma}
\label{lemma-factor-almost-isomorphism}
Let $k$ be an algebraically closed field. Let $f : X' \to X$ be a
finite morphism of algebraic $k$-schemes such that
$\mathcal{O}_X \subset f_*\mathcal{O}_{X'}$ and such that $f$ is an
isomorphism away from a finite set of points. Then there is a factorization
$$
X' = X_n \to X_{n - 1} \to \ldots \to X_1 \to X_0 = X
$$
such that each $X_i \to X_{i - 1}$ is either the glueing of
two points or the squishing of a tangent vector
(see Examples \ref{example-glue-points} and
\ref{example-squish-tangent-vector}).
\end{lemma}

\begin{proof}
Let $U \subset X$ be the maximal open set over which $f$ is an isomorphism.
Then $X \setminus U = \{x_1, \ldots, x_n\}$ with $x_i \in X(k)$.
We will consider factorizations $X' \to Y \to X$ of $f$ such that
both morphisms are finite and
$$
\mathcal{O}_X \subset g_*\mathcal{O}_Y \subset f_*\mathcal{O}_{X'}
$$
where $g : Y \to X$ is the given morphism. By assumption
$\mathcal{O}_{X, x} \to (f_*\mathcal{O}_{X'})_x$ is an isomorphism
onless $x = x_i$ for some $i$. Hence the cokernel
$$
f_*\mathcal{O}_{X'}/\mathcal{O}_X = \bigoplus \mathcal{Q}_i
$$
is a direct sum of skyscraper sheaves $\mathcal{Q}_i$ supported at
$x_1, \ldots, x_n$.
Because the displayed quotient is a coherent $\mathcal{O}_X$-module,
we conclude that $\mathcal{Q}_i$ has finite length over
$\mathcal{O}_{X, x_i}$. Hence we can argue
by induction on the sum of these lengths, i.e., the length of
the whole cokernel.

\medskip\noindent
If $n > 1$, then we can define an $\mathcal{O}_X$-subalgebra
$\mathcal{A} \subset f_*\mathcal{O}_{X'}$ by taking the inverse
image of $\mathcal{Q}_1$. This will give a nontrivial factorization
and we win by induction.

\medskip\noindent
Assume $n = 1$. We abbreviate $x = x_1$. Consider the finite
$k$-algebra extension
$$
A = \mathcal{O}_{X, x} \subset (f_*\mathcal{O}_{X'})_x = B
$$
Note that $\mathcal{Q} = \mathcal{Q}_1$ is the skyscraper sheaf
with value $B/A$.
We have a $k$-subalgebra $A \subset A + \mathfrak m_A B \subset B$.
If both inclusions are strict, then we obtain a nontrivial
factorization and we win by induction as above.
If $A + \mathfrak m_A B = B$, then $A = B$ by Nakayama, then
$f$ is an isomorphism and there is nothing to prove.
We conclude that we may assume $B = A + \mathfrak m_A B$.
Set $C = B/\mathfrak m_A B$. If $C$ has more than $2$
$k$-subalgebras, then we obtain a subalgebra between $A$
and $B$ by taking the inverse image in $B$. Thus we may
assume $C$ has exactly $2$ $k$-subalgebras. Thus $C = k \times k$
or $C = k[\epsilon]$ by Lemma \ref{lemma-no-in-between-over-k}.
In this case $f$ is correspondingly the glueing two points or the
squishing of a tangent vector.
\end{proof}